\section{系统总体架构}
\label{sec:system}

\subsection{PS--PL 边界}

\cref{fig:system} 给出本文系统边界。KV260 搭载 XCK26 SoC，处理系统侧为
四核 Cortex-A53，PL 侧部署 \LASA 和四路 AXI DMA。DDR 同时容纳常驻参数、
量化输入、两个 HWC ping-pong 区、route skip、分支临时张量和双检测头。
四路 DMA 分别搬运 bias、weight、IFM 和 OFM；控制面使用 AXI-Lite 设置层
描述符、量化参数、LUT、地址、字节数和启动命令，并读取完成状态、错误位
与性能计数器。DMA 数据口和共享 DDR 通过 PS 高性能端口连接，PL 不直接
承担文件系统、TCP 或 COCO 数据集协议。

\begin{figure*}[t]
\centering
\begin{tikzpicture}[>=Latex,node distance=7mm and 10mm,
  box/.style={draw,rounded corners,minimum height=10mm,minimum width=28mm,align=center,fill=LasaLight},
  cpu/.style={box,fill=LasaOrange!18}, accel/.style={box,fill=LasaBlue!15},
  mem/.style={box,fill=LasaGreen!15}, io/.style={box,fill=gray!12},
  flow/.style={->,thick,LasaGray}]
\node[io] (eth) {千兆以太网/SD\\输入与结果归档};
\node[cpu,right=of eth,minimum width=37mm] (a53) {四核 Cortex-A53\\DAG调度/张量算子\\双头 decode/NMS};
\node[mem,right=of a53,minimum width=34mm] (ddr) {DDR tensor arena\\参数常驻/双尺度分支};
\node[accel,below=15mm of ddr,minimum width=40mm] (lasa) {LASA PL IP\\13 个 INT8 卷积};
\node[accel,left=of lasa,minimum width=35mm] (dma) {4 路 AXI DMA\\bias/weight/IFM/OFM};
\node[io,left=of dma] (ctrl) {AXI-Lite\\配置/计数器};
\draw[flow] (eth)--(a53);
\draw[flow,<->] (a53)--(ddr);
\draw[flow,<->] (ddr)--node[right,font=\scriptsize]{HP0--HP3}(lasa);
\draw[flow,<->] (dma)--(lasa);
\draw[flow] (ctrl)--(lasa);
\draw[flow] (a53) |- (ctrl);
\node[draw,dashed,fit=(ddr)(lasa)(dma)(ctrl),inner sep=5mm,label={[font=\small]below:XCK26 PL/PS 共享系统}] {};
\end{tikzpicture}
\caption{端到端 MPSoC 系统边界。网络/SD 只负责数据供给与证据封存；
resident 计时从量化输入已驻 DDR 开始，到双头检测结果完成结束。}
\label{fig:system}
\end{figure*}


完整推理采用“PL 规则密集算子、A53 图和不规则算子”的边界。13 个卷积均由
\LASA 执行，包括 $C_{out}=255$ 的 P4/P5 检测头。A53 执行普通 2$\times$2
池化、m10 后的右/下量化零点 padding 加 stride-1 pool、nearest 上采样、
route 分支重量化与 HWC concat，以及双头 decode/NMS。该划分并非声称
非卷积永久适合 CPU，而是为了在不改变 r5 bitstream 的前提下验证完整网络；
它同时使本文可以单独测量 PL 卷积和 A53 辅助算子的代价。

\subsection{常驻参数与 tensor arena}

参数包使用独立的 bias 与 weight payload，并由 manifest 描述 13 层的偏移、
长度、tile 重复、量化域和哈希。当前安全 tile 表下，权重约 18.614 MB、
bias 64,256 B，分别位于 20 MiB 和 64 KiB 的既有 DDR 窗口内。软件启动时
只传输和验证一次参数，之后每张图只更新输入与 feature tensor；因此
resident 口径排除参数装载。

feature 生命周期不是简单的单链 ping-pong。m8 的
26$\times$26$\times$256 输出必须保留为 route skip，同时 pool9 产生主干
13$\times$13$\times$256；m14 同时馈入 m15 和 m16；m16 上采样后与 m8
经过重量化的 skip 拼接为 26$\times$26$\times$384。软件以静态 liveness
表给每条边分配 64 B 对齐的低 2 GiB DDR 地址，所有 DMA 可见地址均小于
\code{0x80000000}。cache 操作遵循“CPU 写后 flush、DMA 写后 invalidate”，
协议同时校验字节数、TKEEP/TLAST、完成状态和计数器，防止因 cache 命中
掩盖 DMA 不一致。

\subsection{输入、输出与持久执行}

系统提供 SD 冷启动和以太网持久 runner 两条通道。SD 卡包含 BOOT.bin、
启动脚本、参数和可选数据集，适合无主机冷启动与稳定性验证；以太网通道
使用静态 IP 和 chunk 级 CRC/sequence 协议，适合 5000 张数据集和性能测试。
二者最终调用同一层调度、tensor op 和 decode 核心，因此传输通道不是模型
实现分叉。

以太网会话一次校验 BIT/XSA/ELF/参数哈希，输入按固定 416$\times$416 HWC
量化包发送。板端可以返回三种逻辑输出：双 raw head、A53 detections 或仅
timing/CRC。raw-head 模式用于隔离网络整数语义；product 模式用于验证 A53
decode/NMS；performance 模式关闭 raw dump 和计时区内 UART/JSON。系统由此
区分四种时间：
\begin{itemize}[leftmargin=1.8em]
  \item \textbf{resident}：输入张量已在 DDR，到 detections 就绪；
  \item \textbf{PL}：13 个卷积 dispatch 的板端周期累计；
  \item \textbf{pipeline}：主机发送、板端计算和结果返回的 chunk wall time；
  \item \textbf{cold boot/SD I/O}：供部署重复性报告，不与 resident 混合。
\end{itemize}

\subsection{完整网络执行次序}

\cref{fig:dag} 展示 13 个 PL 卷积和 A53 tensor op 的数据依赖。m8 以 pool
bypass 输出 route；A53 完成 pool9 后继续 m10。m10 的 13$\times$13$\times$512
张量先在右侧和底部补一个量化零点像素，再做 2$\times$2/stride-1 maxpool，
与原网络的非对称 padding 语义一致。m14 输出分为 P5 分支和 P4 neck：前者
经 m15 和 P5 detect，后者经 m16、nearest 上采样、route requant-concat、
m19 和 P4 detect。双 raw head 分别为 172,380 B 和 43,095 B，最后由同一
accuracy 或 demo profile 解码。

\begin{figure*}[t]
\centering
\resizebox{0.98\textwidth}{!}{%
\begin{tikzpicture}[>=Latex,node distance=4mm and 5mm,
  conv/.style={draw,rounded corners,fill=LasaBlue!13,minimum width=15mm,minimum height=7mm,font=\scriptsize,align=center},
  cpu/.style={draw,rounded corners,fill=LasaOrange!16,minimum width=15mm,minimum height=7mm,font=\scriptsize,align=center},
  head/.style={draw,rounded corners,fill=LasaGreen!18,minimum width=18mm,minimum height=7mm,font=\scriptsize,align=center},
  f/.style={->,thick}]
\node[conv] (m0) {m0}; \node[cpu,right=of m0] (p1) {pool1};
\node[conv,right=of p1] (m2) {m2}; \node[cpu,right=of m2] (p3) {pool3};
\node[conv,right=of p3] (m4) {m4}; \node[cpu,right=of m4] (p5) {pool5};
\node[conv,right=of p5] (m6) {m6}; \node[cpu,right=of m6] (p7) {pool7};
\node[conv,below=8mm of m6] (m8) {m8}; \node[cpu,right=of m8] (p9) {pool9};
\node[conv,right=of p9] (m10) {m10}; \node[cpu,right=of m10] (p12) {pad+pool12};
\node[conv,right=of p12] (m13) {m13}; \node[conv,right=of m13] (m14) {m14};
\node[conv,below=8mm of m13] (m15) {m15}; \node[head,right=of m15] (p5h) {P5 head\\13$\times$13$\times$255};
\node[conv,below=8mm of m14] (m16) {m16}; \node[cpu,right=of m16] (up) {nearest$\times$2};
\node[cpu,right=of up] (cat) {requant\\concat}; \node[conv,right=of cat] (m19) {m19};
\node[head,right=of m19] (p4h) {P4 head\\26$\times$26$\times$255};
\foreach \a/\b in {m0/p1,p1/m2,m2/p3,p3/m4,m4/p5,p5/m6,m6/p7,p7/m8,m8/p9,p9/m10,m10/p12,p12/m13,m13/m14,m14/m15,m15/p5h,m14/m16,m16/up,up/cat,cat/m19,m19/p4h}{\draw[f] (\a)--(\b);}
\draw[f,LasaRed] (m8.south) |- (cat.south);
\node[draw,dashed,fit=(m0)(m2)(m4)(m6)(m8)(m10)(m13)(m14)(m15)(m16)(m19)(p4h)(p5h),inner sep=3mm,label={[font=\scriptsize]below:蓝/绿：13 次 PL 卷积 dispatch}] {};
\node[font=\scriptsize,LasaRed,anchor=north] at (m8.south) {route skip};
\end{tikzpicture}
}
\caption{完整 COCO80 YOLOv3-tiny DAG。m8 的 26$\times$26$\times$256
输出既形成 route skip，也经 pool9 进入低分辨率主干；两头输出共 2535 个 anchor。}
\label{fig:dag}
\end{figure*}


这种系统边界保留了两个重要性质。第一，所有卷积都经过同一硬件数值语义，
因此可用一个 RTL-semantic host 模型逐层比较；第二，训练或 QAT 只替换参数
manifest，不改变 13 层图和硬件调度合同。当前 epoch-1 参数精度不足是模型
release 问题，而不是通过重新定义板端后处理可以规避的硬件问题。
